% ── Chapter 2: Related Work and Background Theory ─────────────
\chapter{งานวิจัยที่เกี่ยวข้องและทฤษฎีพื้นฐาน (Related Work and Background Theory)}

\section{การทดสอบความไวต่อยาปฏิชีวนะ (Antibiotic Susceptibility Testing)}

วิธีการทดสอบ Kirby-Bauer disk diffusion ได้รับการอธิบายและกำหนดมาตรฐานอย่างเป็นระบบครั้งแรกโดย Bauer, Kirby, Sherris และ Turck ในปี 1966 นับแต่นั้นมา วิธีนี้ได้กลายเป็นเทคนิคที่ได้รับการยอมรับอย่างแพร่หลายที่สุดในการทดสอบความไวต่อยาต้านจุลชีพในห้องปฏิบัติการจุลชีววิทยาคลินิกทั่วโลก \cite{asm2009kirby,clsi2025} ความนิยมอย่างต่อเนื่องของวิธีนี้มาจากต้นทุนอุปกรณ์ที่ต่ำ ขั้นตอนที่เรียบง่าย และความสามารถในการประยุกต์ใช้กับเชื้อก่อโรคแบคทีเรียที่เจริญเติบโตในสภาวะที่มีออกซิเจนได้เกือบทุกชนิด หลักการทำงานของวิธีนี้อาศัยการแพร่กระจายของยาปฏิชีวนะผ่านเนื้อวุ้นกึ่งแข็ง กล่าวคือ เมื่อนำแผ่นดิสก์กระดาษที่อาบน้ำยาปฏิชีวนะไปวางบนผิววุ้นที่เพาะเชื้อไว้ โมเลกุลของยาจะแพร่กระจายออกไปเป็นรัศมีรอบแผ่นดิสก์ ทำให้เกิดระดับความเข้มข้นที่ลดลงแบบเอ็กซ์โพเนนเชียลตามระยะทาง บริเวณที่ความเข้มข้นของยาปฏิชีวนะสูงกว่าค่าความเข้มข้นต่ำสุดที่ยับยั้งเชื้อได้ (Minimum Inhibitory Concentration: MIC) การเจริญเติบโตของแบคทีเรียจะถูกยับยั้ง ปรากฏเป็นวงใสทรงกลมบนพื้นผิววุ้นที่เต็มไปด้วยแบคทีเรีย

ขนาดเส้นผ่านศูนย์กลางของบริเวณยับยั้งเชื้อ (Zone of Inhibition: ZOI) มีความสัมพันธ์แบบผกผันกับค่า MIC โดยเชื้อที่มีความไวต่อยาสูง (ค่า MIC ต่ำ) จะเกิด ZOI ขนาดใหญ่ ในขณะที่เชื้อที่ดื้อยาตามธรรมชาติหรือดื้อยาที่ได้รับมาจะเกิดโซนขนาดเล็กหรือไม่มีเลย ความสัมพันธ์ระหว่างขนาดเส้นผ่านศูนย์กลาง ZOI และ MIC ไม่ได้เป็นแบบเชิงเส้น แต่เป็นไปตามสมการทางคณิตศาสตร์ที่ขึ้นอยู่กับปัจจัยต่างๆ เช่น ความหนาของชั้นวุ้น ปริมาณยาบนแผ่นดิสก์ สัมประสิทธิ์การแพร่ และระยะเวลาในการบ่มเชื้อ ซึ่งปัจจัยทั้งหมดนี้ถูกควบคุมอย่างเข้มงวดตามเกณฑ์มาตรฐาน ในการใช้งานทางคลินิก ขนาดเส้นผ่านศูนย์กลาง ZOI ที่วัดได้จะถูกนำไปเทียบกับเกณฑ์จุดตัด (Interpretive Breakpoints) เพื่อจำแนกความไวต่อยาออกเป็นกลุ่ม (Susceptible, Intermediate หรือ Resistant) โดยเกณฑ์จุดตัดเหล่านี้ได้มาจากการรวบรวมข้อมูลทางคลินิกและจุลชีววิทยา วิธีการนี้ถูกนำไปใช้อย่างกว้างขวางกับแบคทีเรียที่มีความสำคัญทางคลินิกหลายชนิด รวมถึง \textit{Bacillus cereus} \cite{fitwi2020} และ \textit{Bacillus subtilis} \cite{jiang2023} ตลอดจนเชื้อก่อโรคอื่นๆ ที่พบประจำในห้องปฏิบัติการ

\subsection{มาตรฐานทางคลินิก: CLSI และ EUCAST (Clinical Standards: CLSI and EUCAST)}

เกณฑ์จุดตัด (Interpretive Breakpoints) ซึ่งเป็นเกณฑ์ขนาดเส้นผ่านศูนย์กลางที่ใช้จำแนกเชื้อแบคทีเรียว่าเป็น ไวต่อยา (S), กึ่งไว (I) หรือ ดื้อยา (R) นั้น ได้รับการกำหนดและปรับปรุงอย่างสม่ำเสมอโดยองค์กรระดับนานาชาติ 2 แห่ง ได้แก่ Clinical and Laboratory Standards Institute (CLSI) ในสหรัฐอเมริกา และ European Committee on Antimicrobial Susceptibility Testing (EUCAST) ในยุโรป \cite{eucast2023} แม้ว่าทั้งสององค์กรจะมีเป้าหมายทางคลินิกที่ตรงกัน แต่ระเบียบวิธีในการพิจารณาเกณฑ์จุดตัดกลับมีความแตกต่างกันในบางประเด็น ส่งผลให้เกณฑ์ที่ประกาศใช้สำหรับยาและเชื้อแบคทีเรียคู่เดียวกันอาจคลาดเคลื่อนกันได้หลายมิลลิเมตร

CLSI ได้ประกาศใช้มาตรฐานการทดสอบ Disk Diffusion ในเอกสาร M02 ซึ่งปัจจุบันปรับปรุงถึงฉบับที่ 14 \cite{clsi2025} ในขณะที่เกณฑ์ประสิทธิภาพการทดสอบความไวต่อยาระบุอยู่ในเอกสาร M100 ฉบับที่ 35 \cite{clsi_m100_2025} เกณฑ์จุดตัดเหล่านี้กำหนดขึ้นโดยคณะทำงานสหสาขาวิชาชีพ ซึ่งรวบรวมข้อมูลการกระจายตัวของค่า MIC จากเชื้อที่แยกได้ทางคลินิกจำนวนมาก ผสานกับการจำลองทางเภสัชจลนศาสตร์และเภสัชพลศาสตร์ (PK/PD) ของการให้ยาในขนาดมาตรฐาน และข้อมูลผลลัพธ์ทางคลินิกจากการทดลองแบบสุ่มที่มีกลุ่มควบคุม สำหรับหมวดหมู่กึ่งไว (Intermediate: I) ในระบบ CLSI เดิมทีหมายถึงช่วงกันชน (Buffer Zone) เพื่อรองรับความคลาดเคลื่อนทางเทคนิคที่อาจทำให้แปลผลผิดพลาด ทว่าในฉบับหลังๆ ได้มีการปรับความหมายให้สอดคล้องกับแนวคิดของ EUCAST มากขึ้น คือเป็นหมวดหมู่ ``Susceptible, Increased Exposure'' ซึ่งหมายถึงเชื้อจะไวต่อยาก็ต่อเมื่อมีการเพิ่มขนาดยาให้ผู้ป่วยได้รับยาสูงสุดตามความเหมาะสม

ทางด้าน EUCAST จะมีการอัปเดตตารางเกณฑ์จุดตัดเป็นประจำทุกปี (เวอร์ชันที่ใช้อ้างอิงในโครงงานนี้คือ 13.1 \cite{eucast2023}) โดยพิจารณาเกณฑ์จากเป้าหมาย PK/PD, ค่าจุดตัดทางระบาดวิทยา (Epidemiological Cutoff Values: ECOFFs) และข้อมูลทางคลินิก \cite{giske2022} ค่า ECOFFs คือขีดจำกัดบนของการกระจายตัวของค่า MIC ในประชากรแบคทีเรียสายพันธุ์ดั้งเดิม (Wild-type) ที่ยังไม่มีกลไกการดื้อยา ดังนั้นเชื้อแบคทีเรียใดที่มีค่า MIC สูงกว่า ECOFF จะถูกจัดเป็นกลุ่มดื้อยา (Non-wild-type: NWT) และถือว่ามีกลไกการดื้อยาเกิดขึ้น เกณฑ์จุดตัดของ EUCAST มักจะเข้มงวดกว่าในการแยกกลุ่มเชื้อที่ไวต่อยาออกจากกลุ่มที่ดื้อยา ซึ่งสะท้อนถึงกรอบแนวคิดทางเภสัชวิทยาของยุโรปที่แตกต่างไปจากของสหรัฐอเมริกา \cite{kahlmeter2019} ระบบที่พัฒนาขึ้นในโครงงานนี้ได้รวบรวมเกณฑ์จุดตัดจากทั้งสองมาตรฐานไว้ในฐานข้อมูล เพื่อให้ห้องปฏิบัติการสามารถเลือกใช้มาตรฐานที่สอดคล้องกับระเบียบข้อบังคับและบริบททางคลินิกของตนเองได้อย่างยืดหยุ่น

\begin{figure}[H]
  \centering
  \safeimg[width=0.78\textwidth]{figures/_review/clsi_breakpoint_table.png}{CLSI breakpoint table for S. aureus}
  \caption{ตัวอย่างตารางเกณฑ์การแปลผลเส้นผ่านศูนย์กลางโซน (Zone diameter breakpoints) สำหรับ
    \textit{Staphylococcus aureus} ตามมาตรฐาน CLSI แสดงช่วง Resistant, Intermediate และ Susceptible
    สำหรับยาปฏิชีวนะแต่ละชนิด \cite{clsi2025}}
  \label{fig:clsi_table}
\end{figure}

\subsection{ความสำคัญทางคลินิกของความแม่นยำในการวัด (Clinical Significance of Measurement Precision)}

ความผิดพลาดในการวัดขนาดเส้นผ่านศูนย์กลางในวิธี Disk Diffusion อาจส่งผลกระทบทางคลินิกที่ร้ายแรงโดยตรง เนื่องจากเกณฑ์จุดตัดถูกกำหนดเป็นตัวเลขจำนวนเต็มในหน่วยมิลลิเมตร ความคลาดเคลื่อนเพียง 1 มิลลิเมตรอาจทำให้ผลการแปลผลเปลี่ยนจากไวต่อยาเป็นดื้อยา หรือในทางกลับกันได้ การจำแนกผลผิดพลาดนี้ก่อให้เกิดอันตราย 2 รูปแบบหลัก ได้แก่ ผลไวต่อยาปลอม (False Susceptibility) คือการระบุว่าเชื้อดื้อยากลับเป็นเชื้อที่ไวต่อยา ซึ่งอาจนำไปสู่ความล้มเหลวในการรักษาและเป็นอันตรายถึงชีวิต และผลดื้อยาปลอม (False Resistance) คือการระบุว่าเชื้อที่ไวต่อยากลับเป็นเชื้อดื้อยา ซึ่งจะบีบให้แพทย์ต้องเปลี่ยนไปใช้ยาทางเลือกที่สองหรือสามโดยไม่จำเป็น เป็นการเพิ่มทั้งค่าใช้จ่ายและเร่งให้เกิดปัญหาเชื้อดื้อยาในอนาคต Humphries และคณะ \cite{humphries2018} ได้กำหนดให้ค่าความคลาดเคลื่อนสัมบูรณ์เฉลี่ย (MAE) ที่น้อยกว่า 3 มม. เป็นเกณฑ์ความแม่นยำที่ยอมรับได้ทางคลินิกสำหรับเครื่องมือวัด AST อัตโนมัติ ซึ่งโครงงานนี้ได้นำเกณฑ์ดังกล่าวมาเป็นเป้าหมายหลักในการประเมินประสิทธิภาพ

\section{Deep Learning สำหรับการตรวจจับวัตถุ (Deep Learning for Object Detection)}

การตรวจจับวัตถุ (Object Detection) ในสาขาคอมพิวเตอร์วิทัศน์ (Computer Vision) คืองานที่มุ่งเน้นการระบุตำแหน่งและจำแนกประเภทของวัตถุหนึ่งชิ้นหรือมากกว่าที่ปรากฏอยู่ในภาพเดียวกัน ในมุมมองกว้างๆ การเรียนรู้ของเครื่อง (Machine Learning) ประกอบด้วยกระบวนทัศน์หลายรูปแบบ เช่น การเรียนรู้แบบมีผู้สอน (Supervised Learning), แบบไม่มีผู้สอน (Unsupervised Learning) และแบบเสริมกำลัง (Reinforcement Learning) \cite{alam2023ml,murphy2025rl} ทว่าสำหรับงานด้านการจดจำภาพ (Image Recognition) แล้ว Supervised Deep Learning พิสูจน์แล้วว่ามีประสิทธิภาพสูงสุด โดยมี Convolutional Neural Networks (CNNs) เป็นรากฐานของโมเดลตรวจจับวัตถุส่วนใหญ่ในปัจจุบัน \cite{purwono2023cnn} งานทบทวนวรรณกรรมอย่างเป็นระบบด้าน Deep Learning ในทางการแพทย์ยังยืนยันความเหมาะสมของเทคโนโลยีนี้กับงานวิเคราะห์ภาพทางคลินิกในหลายบริบท \cite{egger2022} ระบบตรวจจับวัตถุจะต้องส่งออกข้อมูลสองส่วนสำหรับวัตถุแต่ละชิ้นที่พบ นั่นคือ กรอบล้อมรอบ (Bounding box) ที่ระบุขอบเขตทางพื้นที่ของวัตถุ และป้ายกำกับคลาส (Class label) ที่ระบุประเภทของวัตถุ งานนี้จึงมีความซับซ้อนกว่าการจำแนกประเภทภาพ (Image Classification) ทั่วไปที่มักกำหนดป้ายกำกับเพียงอันเดียวให้กับทั้งภาพ เนื่องจากระบบจะไม่ทราบล่วงหน้าว่ามีวัตถุจำนวนเท่าใดและอยู่ตรงตำแหน่งใดบ้าง แต่ต้องประมวลผลสรุปเอาจากข้อมูลในภาพเอง ทั้งนี้ แนวทาง Deep Learning ได้เข้ามามีบทบาทนำในการสร้างมาตรฐานใหม่ให้กับการตรวจจับวัตถุ นับตั้งแต่การมาถึงของกลุ่มโมเดล Region-based CNN (R-CNN) และตระกูล YOLO ซึ่งต่างก็ได้รับการพัฒนาอย่างต่อเนื่องทั้งในด้านความแม่นยำและความเร็วในการประมวลผล

ในปัจจุบัน สถาปัตยกรรมของโมเดลตรวจจับวัตถุแบบ Deep Learning มักถูกแบ่งออกเป็น 2 กลุ่มหลัก ได้แก่ ตัวตรวจจับแบบสองขั้นตอน (Two-stage Detectors) ซึ่งมีกระบวนการทำงานโดยขั้นแรกจะสร้างบริเวณที่น่าจะเป็นวัตถุ (Region Proposals) ขึ้นมาก่อน แล้วจึงนำแต่ละบริเวณไปจำแนกประเภทในขั้นตอนต่อไป และ ตัวตรวจจับแบบขั้นตอนเดียว (One-stage หรือ Single-shot Detectors) ซึ่งจะทำนาย Bounding boxes และค่าความน่าจะเป็นของคลาสออกมาโดยตรงผ่านการประมวลผลภาพทั้งภาพเพียงรอบเดียว (Single pass) แต่ละสถาปัตยกรรมต่างก็มีข้อดีและข้อเสียที่แตกต่างกันไป โดยจะต้องชั่งน้ำหนักระหว่างความแม่นยำในการตรวจจับ ทรัพยากรที่ใช้ในการคำนวณ และความหน่วงเวลา (Latency) ในการประมวลผล

\subsection{Faster R-CNN}

โมเดล Faster R-CNN (Region-based Convolutional Neural Network) ซึ่งถูกนำเสนอโดย Ren และคณะในปี 2015 ถือเป็นก้าวสำคัญที่สามารถรวมกระบวนการสร้างบริเวณที่น่าจะเป็นวัตถุเข้ากับส่วนประมวลผลการตรวจจับ (Detection Head) ให้อยู่ในโครงข่ายเดียวกัน และสามารถทำการฝึกสอนแบบเบ็ดเสร็จ (End-to-end) ได้ \cite{ren2015faster} ก่อนหน้านี้ โมเดลในตระกูลเดียวกันอย่าง R-CNN และ Fast R-CNN จำเป็นต้องใช้อัลกอริทึม Selective Search ที่แยกออกไปต่างหากเพื่อสร้างบริเวณข้อเสนอ (Region Proposals) ซึ่งใช้ทรัพยากรการคำนวณสูงและกลายเป็นคอขวดที่ทำให้ระบบทำงานช้า Faster R-CNN ได้แก้ปัญหานี้โดยการนำเครือข่าย Region Proposal Network (RPN) เข้ามาแทนที่ โดย RPN จะใช้ประโยชน์จาก Convolutional Backbone ร่วมกับส่วนประมวลผลการตรวจจับ ทำให้สามารถสร้างบริเวณข้อเสนอได้ด้วยความเร็วเกือบเรียลไทม์

\begin{figure}[H]
  \centering
  \safeimg[width=0.82\textwidth]{figures/faster_rcnn_architecture.png}{Faster R-CNN two-stage detection architecture}
  \caption{สถาปัตยกรรมของ Faster R-CNN แสดง Convolutional backbone ที่ใช้ร่วมกัน,
    Region Proposal Network (RPN) ที่สร้างบริเวณวัตถุที่เป็นไปได้ และหัวการตรวจจับ
    ที่ทำการจำแนกและปรับแต่งข้อเสนอแต่ละรายการ \cite{ren2015faster}}
  \label{fig:faster_rcnn}
\end{figure}

สถาปัตยกรรมของ Faster R-CNN มีกระบวนการทำงานแบ่งออกเป็น 2 ขั้นตอนหลัก ในขั้นตอนแรก Convolutional Backbone (โครงงานนี้เลือกใช้โมเดล ResNet-50 \cite{he2016deep}) จะทำหน้าที่สกัดแผนผังคุณลักษณะเชิงพื้นที่ (Spatial Feature Maps) ที่มีความซับซ้อนออกจากภาพตั้งต้นแบบเต็มภาพ จากนั้น RPN จะใช้เครือข่ายย่อยกวาดเลื่อน (Slide) ไปบนแผนผังคุณลักษณะเหล่านี้ เพื่อประเมินความน่าจะเป็นที่แต่ละตำแหน่งเชิงพื้นที่จะมีวัตถุอยู่ (Objectness Scores) พร้อมกับคำนวณค่าปรับแก้ตำแหน่ง (Regression Offsets) สำหรับ Bounding Box โดยพิจารณาจากชุดของ Anchor Boxes ที่ถูกกำหนดขนาดและสัดส่วนล่วงหน้า Anchor Boxes ที่มีคะแนนความน่าจะเป็นสูงกว่าเกณฑ์ที่กำหนดจะถูกคัดเลือกไว้เป็นบริเวณข้อเสนอ (โดยทั่วไปจะเหลือประมาณ 256--512 บริเวณต่อภาพ หลังจากผ่านกระบวนการ Non-Maximum Suppression) ต่อมาในขั้นตอนที่สอง ระบบจะใช้ฟังก์ชัน RoI Pooling (หรือ RoI Align) เพื่อสกัดเวกเตอร์คุณลักษณะที่มีขนาดคงที่สำหรับแต่ละบริเวณข้อเสนอที่ได้จากขั้นตอนแรก โดยดึงข้อมูลจาก Backbone Feature Map เดียวกัน ท้ายที่สุด เครือข่าย Fully-Connected (Classification Head) จะทำหน้าที่ระบุคลาสของวัตถุและปรับแต่งขอบเขตของ Bounding Box แต่ละอันให้มีความแม่นยำมากยิ่งขึ้น

การเลือกใช้โมเดล ResNet-50 เป็น Backbone สำหรับสกัดคุณลักษณะมีความเหมาะสมอย่างยิ่งกับงานนี้ งานวิจัยของ He และคณะ แสดงให้เห็นว่า โครงสร้างการเชื่อมต่อแบบข้ามชั้น (Residual Connections) ซึ่งเป็นการส่งต่อข้อมูลข้ามชั้น Convolution 1 ชั้นหรือมากกว่านั้น ช่วยแก้ปัญหากราเดียนต์หายไป (Vanishing Gradient) ทำให้สามารถฝึกสอนโครงข่ายที่มีความลึกมากๆ ได้อย่างมีประสิทธิภาพ \cite{he2016deep} สถาปัตยกรรมของ ResNet-50 ประกอบไปด้วยเลเยอร์ 50 ชั้นที่จัดเรียงเป็นกลุ่มของ Residual Blocks การที่โมเดลได้รับการฝึกสอนเบื้องต้น (Pretrained) ด้วยชุดข้อมูล ImageNet มาก่อน ทำให้ระบบมีความสามารถในการจับคุณลักษณะพื้นฐานของภาพ เช่น เส้นขอบ พื้นผิว และรูปทรงต่างๆ ซึ่งสามารถถ่ายทอดทักษะความรู้ (Transfer Learning) มาประยุกต์ใช้กับภาพ AST ได้เป็นอย่างดี แม้ว่าชุดข้อมูลภาพ AST ที่ใช้สอนจะมีจำนวนจำกัดก็ตาม ดังนั้น Backbone ResNet-50 ในโครงงานนี้จึงเริ่มต้นด้วยค่าน้ำหนักที่ได้จาก ImageNet ก่อนจะนำมาปรับจูนแบบละเอียด (Fine-tune) กับชุดข้อมูล AST แบบครบวงจร

Faster R-CNN ได้รับการยอมรับอย่างกว้างขวางถึงความแม่นยำในการตรวจจับที่โดดเด่น โดยเฉพาะกับภาพที่มีความซับซ้อน วัตถุมีการทับซ้อนกัน หรืออยู่รวมกันอย่างหนาแน่น ซึ่งตรงกับลักษณะของจานเพาะเชื้อ AST ที่ZOI ของแผ่นดิสก์ที่อยู่ใกล้เคียงกันมักจะทับซ้อนกันบางส่วน กระบวนการทำงานแบบ 2 ขั้นตอนนี้ส่งผลดีให้ Classification Head สามารถวิเคราะห์เฉพาะบริเวณข้อเสนอที่มีคุณภาพและผ่านการกรองเชิงพื้นที่มาแล้ว แทนที่จะต้องกวาดวิเคราะห์แบบสุ่มไปทั่วทั้งภาพ จึงช่วยลดโอกาสในการทายผลผิดพลาดว่ามีวัตถุ (False Positive) อย่างไรก็ดี ข้อจำกัดหลักของ Faster R-CNN คือความเร็วในการทำงาน โดยความหน่วงเวลาจากกระบวนการ 2 ขั้นตอนอาจรับได้สำหรับการวิเคราะห์แบบออฟไลน์ (Offline Analysis) แต่จะกลายเป็นข้อจำกัดสำคัญหากต้องการประมวลผลข้อมูลปริมาณมหาศาลแบบเรียลไทม์

\subsection{YOLOv11}

โมเดล YOLO (You Only Look Once) เป็นตัวตรวจจับแบบขั้นตอนเดียว (Single-stage Detector) ที่ปรับกระบวนทัศน์การตรวจจับวัตถุให้เป็นปัญหาการถดถอย (Regression Problem) แบบรวมศูนย์ แทนที่จะแยกขั้นตอนการสร้างบริเวณข้อเสนอและการจำแนกคลาสออกจากกัน โมเดล YOLO สามารถทำนายพิกัด Bounding Box และค่าความน่าจะเป็นของคลาสออกมาได้โดยตรงจากการประมวลผลภาพรวมเพียงรอบเดียว (Single Pass) \cite{wang2024yolov11,khanam2024yolov11} การออกแบบนี้ช่วยตัดขั้นตอนการสร้างบริเวณข้อเสนอที่ยุ่งยากออกไป ทำให้ความเร็วในการประมวลผลเพิ่มขึ้นอย่างก้าวกระโดด แม้ว่าในรุ่นแรกๆ อาจจะต้องแลกมาด้วยความแม่นยำที่ลดลงสำหรับวัตถุขนาดเล็กหรือวัตถุที่อยู่รวมกันหนาแน่นก็ตาม ทว่าจากการศึกษาเปรียบเทียบในปัจจุบันได้ยืนยันแล้วว่าประสิทธิภาพของ YOLOv11 นั้นสามารถเทียบเคียงหรือทำได้ดีกว่ารุ่นก่อนหน้า รวมถึงโมเดล Faster R-CNN ในงานตรวจจับที่หลากหลาย \cite{sharma2024yolo_comparison}

\begin{figure}[H]
  \centering
  \safeimg[width=0.82\textwidth]{figures/yolov11_architecture.png}{YOLOv11 single-stage detection architecture}
  \caption{สถาปัตยกรรมของ YOLOv11 แสดงเครือข่าย Backbone สำหรับสกัดคุณลักษณะ,
    ส่วนคอ (neck) แบบ Feature Pyramid สำหรับการตรวจจับหลายมาตราส่วน
    และหัวการตรวจจับแยกส่วน (Decoupled heads) สำหรับการถดถอย Bounding box และการทำนายคลาส \cite{wang2024yolov11}}
  \label{fig:yolov11}
\end{figure}

YOLOv11 เป็นเวอร์ชันล่าสุดของตระกูล YOLO ในช่วงเวลาที่ดำเนินโครงงานนี้ ซึ่งพัฒนาโดย Ultralytics โดยได้รวบรวมการปรับปรุงทางสถาปัตยกรรมที่สำคัญหลายประการ \cite{wang2024yolov11} ส่วน Backbone เลือกใช้บล็อก C3k2 รูปแบบใหม่ (การออกแบบ Convolution แบบ Reparameterization) เพื่อยกระดับประสิทธิภาพในการสกัดคุณลักษณะ ส่วนคอ (Neck) ใช้เครือข่าย Path Aggregation Feature Pyramid Network (PA-FPN) ในการผสมผสานคุณลักษณะจาก Backbone ในหลายระดับชั้น ทำให้สามารถตรวจจับวัตถุได้หลายขนาดพร้อมกัน ตั้งแต่แผ่นดิสก์ยาปฏิชีวนะที่มีขนาดเล็กไปจนถึง ZOI ขนาดใหญ่ที่อาจกินพื้นที่ส่วนใหญ่ของจานเพาะเชื้อ นอกจากนี้ ส่วนหัวของการตรวจจับ (Detection Head) ยังถูกแยกโครงสร้างออกจากกัน (Decoupled Head) โดยแบ่งสาขาการทำงานสำหรับการทำนาย Bounding Box (Regression) ออกจากการทำนายคลาส (Classification) อย่างชัดเจน ซึ่งมีผลการศึกษาชี้ว่าช่วยปรับปรุงอัตราการลู่เข้าสู่จุดสมดุล (Convergence) และเพิ่มความแม่นยำให้กับตัวตรวจจับแบบขั้นตอนเดียวได้อย่างมีนัยสำคัญ

โมเดล YOLOv11 มีให้เลือกใช้งาน 5 ขนาด (n, s, m, l, x) เพื่อให้ผู้ใช้สามารถปรับสมดุลระหว่างขีดความสามารถของโมเดลและความเร็วในการประมวลผลได้ตามความเหมาะสม สำหรับโครงงานนี้ได้เลือกใช้รุ่น Nano (n) หรือ YOLOv11n โดยให้ความสำคัญกับความเร็วในการทำงาน (Inference Speed) และความเหมาะสมในการนำไปติดตั้งบนฮาร์ดแวร์ที่มีทรัพยากรจำกัด เช่น คอมพิวเตอร์ในห้องปฏิบัติการที่ไม่มี GPU ประสิทธิภาพสูง หรืออุปกรณ์ปลายทาง (Edge Devices) แม้จะเป็นรุ่นที่เล็กที่สุด แต่ YOLOv11n ก็ยังให้ความแม่นยำในการตรวจจับที่โดดเด่นบนชุดข้อมูล AST ของโครงงาน เนื่องจากได้รับประโยชน์จากการฝึกสอนเบื้องต้น (Pretraining) บนชุดข้อมูลมาตรฐานขนาดใหญ่อย่าง COCO ก่อนที่จะนำมาปรับจูน (Fine-tune) กับภาพ AST ที่ได้ทำป้ายกำกับ (Label) ไว้ โมเดลนี้ถูกกำหนดให้ตรวจจับวัตถุ 2 คลาส ได้แก่ \texttt{antibiotic} (แผ่นดิสก์ยาปฏิชีวนะ ซึ่งใช้เป็นเป้าหมายสำหรับ OCR และใช้อ้างอิงในการปรับเทียบสัดส่วนภาพ) และ \texttt{disk-zone} (ZOI รอบๆ แผ่นยา ซึ่งใช้สำหรับวัดขนาดเส้นผ่านศูนย์กลาง)

\subsection{Hough Circle Transform}

Hough Circle Transform เป็นเทคนิค Computer Vision แบบดั้งเดิมที่ใช้สำหรับการตรวจจับโครงสร้างรูปวงกลมในภาพ โดยไม่ต้องอาศัยโมเดลที่ผ่านการฝึกสอนหรือชุดข้อมูลที่มีป้ายกำกับ วิธีการทำงานเริ่มต้นจากการแปลงภาพจากโดเมนเชิงพื้นที่ (Spatial Domain) ให้อยู่ในรูปแบบของปริภูมิพารามิเตอร์ (Parameter Space) ที่ประกอบด้วย 3 มิติ ได้แก่ พิกัด $x$ และ $y$ ของจุดศูนย์กลางวงกลม และรัศมี $r$ สำหรับทุกๆ จุดที่เป็นเส้นขอบ (Edge Points) ในภาพ (มักหาโดยวิธี Canny Edge Detection) อัลกอริทึมจะทำการ "โหวต" (Vote) ให้กับค่าพารามิเตอร์ของวงกลมที่เป็นไปได้ทั้งหมดที่ลากผ่านจุดขอบนั้น กลุ่มพารามิเตอร์ที่ได้รับคะแนนโหวตสะสมสูงเกินเกณฑ์ที่กำหนด จะถูกระบุว่าเป็นวงกลมที่ตรวจพบในภาพ นอกจากวิธีแบบดั้งเดิมนี้แล้ว ยังมีงานวิจัยที่เสนอแนวทางการตรวจจับวงกลมแบบอื่น เช่น การใช้ Learning Automata \cite{cuevas2014}, Deep Learning \cite{ercan2020} และวิธีการที่อิงกับค่าความชัน (Gradient-based) \cite{shehata2015hough} เพื่อลดข้อจำกัดและเพิ่มความยืดหยุ่นในการใช้งานกับภาพถ่ายจากสภาพแวดล้อมจริง

\begin{figure}[H]
  \centering
  \safeimg[width=0.75\textwidth]{figures/hough_circle_diagram.png}{Hough Circle Transform parameter space voting diagram}
  \caption{ภาพประกอบ Hough Circle Transform จุดขอบแต่ละจุดจะลงคะแนนให้กับกรวยของ
    วงกลมที่เป็นไปได้ในพื้นที่พารามิเตอร์ $(x_c, y_c, r)$
    โดยวงกลมจริงจะสร้างยอดที่โดดเด่นในตัวสะสมคะแนน \cite{shehata2015hough}}
  \label{fig:hough}
\end{figure}

แม้ว่า Hough Circle Transform จะประหยัดทรัพยากรการคำนวณและลดการเกิดการตรวจจับที่ผิดพลาดในภาพที่มีโซนวงกลมที่ชัดเจนภายใต้สภาวะแสงที่ควบคุมได้ดี แต่อัลกอริทึมนี้ก็มีความอ่อนไหวสูงต่อความคลาดเคลื่อนที่พบได้บ่อยในภาพ AST ที่ถ่ายในสภาพจริง ตัวอย่างเช่น เงาสะท้อน แสงที่ไม่สม่ำเสมอ หรือรอยตำหนิบนพื้นผิววุ้น มักจะทำให้เกิดเส้นขอบปลอม ซึ่งไปรบกวนคะแนนสะสมในระบบโหวต นอกจากนี้ โซนที่ไม่เป็นวงกลมสมบูรณ์อันเนื่องมาจากการทับซ้อนกัน ความไม่สม่ำเสมอของวุ้น หรือลักษณะการเจริญเติบโตของแบคทีเรีย จะทำให้คะแนนโหวตกระจายตัวจนไม่เกิดจุดยอดที่ชัดเจน ส่งผลให้ตรวจจับพลาดหรือกะประมาณรัศมีคลาดเคลื่อน ยิ่งไปกว่านั้น วิธีการนี้ยังต้องอาศัยการปรับแต่งค่าพารามิเตอร์ (Tuning) ด้วยมืออย่างระมัดระวัง เช่น สัดส่วนความละเอียดของ Accumulator (\texttt{dp}), ระยะห่างขั้นต่ำระหว่างจุดศูนย์กลาง (\texttt{minDist}) และเกณฑ์ในการหาเส้นขอบ Canny (\texttt{param1}, \texttt{param2}) พารามิเตอร์เหล่านี้มักจะมีความสัมพันธ์กันอย่างซับซ้อน และค่าที่เหมาะสมที่สุดมักจะเปลี่ยนไปตามแต่ละภาพที่มีแสง มุมมอง หรือชนิดแบคทีเรียที่แตกต่างกัน ในโครงงานนี้ การตั้งค่าพารามิเตอร์ของ Hough Circle Transform ทำโดยการค้นหาแบบตาราง (Grid Search) ครอบคลุมทั้งชุดข้อมูล เพื่อให้ได้ประสิทธิภาพสูงสุดที่เป็นไปได้สำหรับเป็นเกณฑ์มาตรฐาน (Baseline) ในการเปรียบเทียบกับโมเดลแบบ Deep Learning

\section{การจดจำอักขระด้วยแสงสำหรับฉลากแผ่นดิสก์ยาปฏิชีวนะ (Optical Character Recognition for Antibiotic Disk Labeling)}

แผ่นดิสก์ยาปฏิชีวนะแต่ละแผ่นจะมีรหัสตัวอักษรและตัวเลข (Alphanumeric Code) สั้นๆ พิมพ์อยู่บนพื้นผิว เพื่อระบุชนิดและความเข้มข้นของตัวยา (เช่น ``OX 1'' หมายถึง Oxacillin 1\,$\mu$g หรือ ``AMX 25'' หมายถึง Amoxicillin 25\,$\mu$g) เทคโนโลยีการจดจำอักขระด้วยแสง (Optical Character Recognition: OCR) ซึ่งได้รับการพัฒนาอย่างต่อเนื่องมาหลายทศวรรษ \cite{wang2023ocr_history} ได้เข้ามามีบทบาทสำคัญในการดึงข้อมูลข้อความเหล่านี้จากภาพถ่ายโดยอัตโนมัติ การอ่านรหัสยาบนแผ่นดิสก์ได้อย่างถูกต้องมีความสำคัญอย่างยิ่งยวดต่อการเชื่อมโยงข้อมูลระหว่าง ZOI ที่ตรวจพบกับชนิดของยาปฏิชีวนะ เพื่อให้ระบบสามารถดึงเกณฑ์จุดตัด (Breakpoints) จากมาตรฐาน CLSI หรือ EUCAST มาแปลผลได้อย่างถูกต้อง อย่างไรก็ตาม การทำ OCR ในบริบทนี้มีความท้าทายสูงด้วยปัจจัยหลายประการ ได้แก่ ขนาดของตัวอักษรที่เล็กมาก (ความสูงเพียง 2--4 มม.) ความคมชัดระหว่างสีหมึกและสีพื้นแผ่นดิสก์ที่แตกต่างกันไปตามยี่ห้อผู้ผลิต ความโค้งมนของขอบแผ่นดิสก์ มุมมองในการถ่ายภาพที่อาจบิดเบี้ยว รวมถึงโอกาสที่แผ่นดิสก์บางส่วนจะถูกบดบังหรือเปื้อน

\subsection{Gemini 2.0 Flash: โมเดลภาษาขนาดใหญ่แบบ Multimodal สำหรับ OCR หลัก (Gemini 2.0 Flash: Multimodal LLM as Primary OCR)}

โครงงานนี้เลือกใช้ \textbf{Gemini 2.0 Flash} \cite{google2025gemini} ซึ่งเป็นโมเดลภาษาขนาดใหญ่แบบ Multimodal (Multimodal Large Language Model: MLLM) ที่พัฒนาโดย Google DeepMind เป็น OCR Engine หลักในการระบุรหัสยาปฏิชีวนะ ความสามารถพิเศษของ Gemini 2.0 Flash คือการเข้าใจบริบท (Context Understanding) ควบคู่กับการรับรู้ทางภาพ (Visual Perception) ในคำสั่งเดียว ซึ่งแตกต่างจากระบบ OCR แบบดั้งเดิมที่แยกขั้นตอนการตรวจจับข้อความและการจดจำตัวอักษรออกจากกัน โมเดลสามารถรับภาพสีต้นฉบับของแผ่นดิสก์ยาพร้อมกับ Prompt ที่บอกบริบทได้โดยตรง แล้วส่งกลับมาเป็นรหัสยาในรูปแบบที่ต้องการ

การนำ MLLM มาใช้แก้ปัญหา OCR ในบริบทนี้มีข้อได้เปรียบหลายประการ ได้แก่ (1) ไม่จำเป็นต้องผ่านขั้นตอนการประมวลผลภาพเบื้องต้น (Preprocessing) เพิ่มเติม เนื่องจากโมเดลสามารถรับภาพสีได้โดยตรงและจัดการกับความแปรปรวนของแสงและมุมมองได้โดยอัตโนมัติ (2) โมเดลสามารถใช้ความรู้เกี่ยวกับรูปแบบรหัสยาปฏิชีวนะที่ได้เรียนรู้ในขั้นตอน Pretraining มาช่วยตีความตัวอักษรที่ไม่ชัดเจนได้ และ (3) โมเดลสามารถจัดการกับทิศทางการหมุนของแผ่นดิสก์ได้โดยอัตโนมัติโดยไม่ต้องทดสอบหลายมุม

\subsection{EasyOCR: ระบบ OCR แบบ Local สำหรับ Fallback (EasyOCR: Local Fallback OCR)}

ในกรณีที่ไม่สามารถเรียกใช้ Gemini API ได้ เช่น เมื่อ API Quota หมดหรือเกิดข้อผิดพลาดในการเชื่อมต่อ ระบบจะสลับไปใช้ \textbf{EasyOCR} \cite{jaidedai2024easyocr} ซึ่งเป็นระบบ OCR แบบ Open-source ที่รันได้ภายในเครื่องเซิร์ฟเวอร์โดยไม่ต้องพึ่งพาบริการภายนอก กระบวนการตรวจจับข้อความของ EasyOCR มีรากฐานมาจากโมเดล CRAFT (Character Region Awareness for Text Detection) \cite{baek2019craft} ซึ่งมุ่งเน้นตรวจจับพื้นที่ของตัวอักษรแต่ละตัว (Character-level) และประเมินความเชื่อมโยงระหว่างตัวอักษรที่อยู่ติดกัน เครือข่าย CRAFT จะสร้างแผนผังคะแนนพื้นที่ (Region Score Map) และแผนผังคะแนนความสัมพันธ์ (Affinity Score Map) จากนั้นพื้นที่ตัวอักษรที่ตรวจพบจะถูกส่งต่อไปยังโครงข่าย CRNN (Convolutional Recurrent Neural Network) เพื่อแปลงลำดับภาพตัวอักษรให้เป็นข้อความด้วยกลไก CTC (Connectionist Temporal Classification) Decoder

ในขั้นตอนการเลือก OCR Engine สำรอง ได้มีการทดสอบเปรียบเทียบระหว่าง EasyOCR \cite{jaidedai2024easyocr} และ TyphoonOCR \cite{typhoonocr2025} ผลการทดสอบพบว่า EasyOCR ให้ความแม่นยำสูงกว่าในการอ่านรหัสตัวอักษรภาษาอังกฤษบนแผ่นดิสก์ยา โดยเฉพาะเมื่อใช้ร่วมกับการประมวลผลภาพเบื้องต้น จึงเลือก EasyOCR เป็น Fallback Engine สำหรับระบบ

\section{เทคนิคการเตรียมการประมวลผลภาพ (Image Preprocessing Techniques)}

การเตรียมการประมวลผลภาพเป็นตัวขับเคลื่อนที่สำคัญของทั้งโมเดลการตรวจจับวัตถุและไปป์ไลน์ OCR ซึ่งส่งผลโดยตรงต่ออัตราส่วนสัญญาณต่อสัญญาณรบกวนของคุณลักษณะที่สกัดโดยอัลกอริทึมปลายน้ำ \cite{qi2021image_enhancement} เทคนิคการแบ่งส่วนภาพ \cite{jasim2021segmentation,dhanachandra2015} และวิธีการแปลงเป็นภาพขาวดำ (binarization) ของเอกสาร \cite{yang2024binarization} ให้ข้อมูลสำหรับการเลือกการเตรียมการประมวลผลในโครงงานนี้ เพื่อจัดการกับการเสื่อมโทรมของภาพเฉพาะที่พบในการถ่ายภาพ AST ในห้องปฏิบัติการ

\begin{itemize}
  \item \textbf{Global Thresholding (วิธีของ Otsu):} อัลกอริทึมของ Otsu ทำหน้าที่คำนวณหาค่าความสว่างที่เป็นเกณฑ์ (Threshold) ที่เหมาะสมที่สุดโดยอัตโนมัติ เพื่อใช้แปลงภาพระดับสีเทา (Grayscale) ให้เป็นภาพขาวดำ (Binary Image) โดยอาศัยหลักการลดค่าความแปรปรวนภายในกลุ่ม (Intra-class Variance) ระหว่างพิกเซลที่เป็นตัวอักษร (Foreground) และพื้นหลัง (Background) ให้น้อยที่สุด โครงงานนี้ประยุกต์ใช้วิธีของ Otsu กับภาพเฉพาะส่วนที่ตัดมาเฉพาะแผ่นดิสก์ยา (Cropped Disk Image) ก่อนจะส่งเข้าสู่ระบบ OCR เพื่อแยกตัวอักษรที่พิมพ์ด้วยหมึกออกจากพื้นผิวและลวดลายของกระดาษแผ่นดิสก์ ซึ่งช่วยยกระดับความแม่นยำในการแยกแยะตัวอักษรได้อย่างมาก

  \item \textbf{Contrast Limited Adaptive Histogram Equalization (CLAHE):} ต่างจากการปรับฮิสโตแกรมแบบปกติ (Global Histogram Equalization) ที่ใช้ฟังก์ชันเดียวในการปรับความเปรียบต่าง (Contrast) ของทั้งภาพ เทคนิค CLAHE จะแบ่งภาพออกเป็นส่วนย่อยๆ (Tiles) แล้วทำการปรับฮิสโตแกรมของแต่ละส่วนอย่างอิสระ (โดยสามารถกำหนดขีดจำกัดเพื่อป้องกันไม่ให้ความเปรียบต่างเพิ่มขึ้นมากเกินไปจนเกิด Noise) จากนั้นจึงทำการเกลี่ยรอยต่อระหว่างส่วนย่อยแต่ละส่วนให้กลมกลืนกัน วิธีนี้ช่วยปรับสมดุลความเปรียบต่างแบบเฉพาะที่ (Local Contrast) ได้เป็นอย่างดี โดยเฉพาะในภาพถ่ายจากห้องปฏิบัติการที่มักพบปัญหาความสว่างไม่สม่ำเสมอ เช่น บริเวณกลางจานเพาะเชื้อสว่างกว่าบริเวณขอบเนื่องจากแสงไฟเพดาน ระบบจะนำ CLAHE มาใช้กับภาพแผ่นดิสก์ที่ตัดมา เพื่อเร่งความคมชัดของตัวอักษรที่จางหรือมองเห็นได้ยากก่อนส่งเข้า OCR

  \item \textbf{Gaussian Blur:} เป็นเทคนิคตัวกรองความถี่ต่ำ (Low-pass Filter) ที่นำมาใช้กับภาพถ่ายเต็มจานก่อนจะเข้าสู่กระบวนการตรวจจับเส้นขอบ (Edge Detection สำหรับวิธี Hough Transform) รวมถึงก่อนกระบวนการหาเกณฑ์ภาพ (Thresholding) การทำ Gaussian Blur ช่วยลบสัญญาณรบกวนความถี่สูงที่เกิดจากเซ็นเซอร์กล้อง ซึ่งมักจะทำให้ระบบเข้าใจผิดว่าเป็นเส้นขอบของวัตถุ ขนาดของตัวกรอง (Kernel Size) จะถูกเลือกอย่างระมัดระวังเพื่อให้สามารถลดสัญญาณรบกวนได้โดยที่ยังคงรักษารูปทรงและขอบเขตที่แท้จริงของ ZOI ไว้ได้

  \item \textbf{การดำเนินการทางสัณฐานวิทยา (Morphological Operations - Dilation และ Erosion):} การขยาย (Dilation) คือการเพิ่มความหนาให้กับพิกเซลของวัตถุในภาพ (เช่น ตัวอักษร) ในขณะที่การกร่อน (Erosion) คือการลดความหนาลง เมื่อนำสองกระบวนการนี้มาใช้งานร่วมกัน (เช่น รูปแบบ Opening หรือ Closing) จะสามารถแก้ปัญหาเส้นตัวอักษรที่ขาดหาย หรือลบจุดรบกวนขนาดเล็กออกจากภาพขาวดำได้ ใน Pipeline ของระบบ OCR การเลือกใช้และระดับของ Morphological Operations จะถูกปรับเปลี่ยนแบบไดนามิกให้เหมาะสมกับความหนาแน่นของพิกเซลตัวอักษรบนแผ่นดิสก์ หากพบว่าตัวอักษรเบาบางหรือเส้นบางเกินไป ระบบจะทำการขยาย (Dilation) เพื่อเชื่อมเส้นที่ขาด แต่หากเส้นตัวอักษรหนาทึบจนกลืนติดกัน ระบบก็จะทำการกร่อน (Erosion) เพื่อให้ตัวอักษรแยกออกจากกันชัดเจนขึ้น

  \item \textbf{การแก้ไขทิศทางการหมุนของภาพ (Image Rotation Correction):} ในทางปฏิบัติของห้องปฏิบัติการ การวางแผ่นดิสก์ยาลงบนจานเพาะเชื้อไม่ได้มีการกำหนดทิศทางที่ตายตัว แผ่นยาอาจถูกวางในมุมใดก็ได้ ทำให้ตัวอักษรที่พิมพ์บนแผ่นยาอาจกลับหัวหรือตะแคงเมื่อเทียบกับกรอบของภาพถ่าย การจัดการความท้าทายนี้แตกต่างกันตาม OCR Engine ที่ใช้งาน สำหรับ Gemini 2.0 Flash ซึ่งเป็น Engine หลัก โมเดลจะจัดการทิศทางโดยอัตโนมัติจาก Pretraining โดยไม่ต้องหมุนภาพ ในขณะที่ EasyOCR ซึ่งเป็น Fallback Engine จะรันซ้ำใน \textbf{24 มุมหมุน} (ครอบคลุมตั้งแต่ $-180^\circ$ ถึง $+165^\circ$ โดยเพิ่มทีละ $15^\circ$) โดยระบบจะเลือกผลลัพธ์จากมุมที่โมเดลให้ค่าความเชื่อมั่น (Confidence Score) สูงที่สุด
\end{itemize}

\section{งานวิจัยที่เกี่ยวข้องในการวิเคราะห์ AST อัตโนมัติ (Related Work in Automated AST Analysis)}

\subsection{แนวทาง Deep Learning สำหรับการวิเคราะห์ภาพ AST (Deep Learning Approaches for AST Image Analysis)}

การประยุกต์ใช้เทคโนโลยี Deep Learning และ Computer Vision ในการวิเคราะห์ภาพ AST ถือเป็นสาขาวิจัยที่มีความตื่นตัวอย่างมากในช่วงทศวรรษที่ผ่านมา โดยได้รับแรงผลักดันจากความต้องการเครื่องมือที่ตอบโจทย์ทางคลินิกผนวกกับความก้าวหน้าอย่างก้าวกระโดดของอัลกอริทึม ซึ่งสอดคล้องกับแรงจูงใจในการริเริ่มโครงงานนี้ งานวิจัยของ Bilal และคณะ \cite{bilal2025} ได้ทบทวนบทบาทที่กว้างขึ้นของ AI และ Machine Learning ในการพยากรณ์และต่อสู้กับปัญหาเชื้อดื้อยาอย่างครอบคลุม ในขณะที่ Liao และคณะ \cite{liao2025} ได้เจาะลึกถึงความก้าวหน้าในงานวิจัยด้านการทดสอบความไวต่อยาที่ขับเคลื่อนด้วย AI โดยเฉพาะ
ส่วน Pascucci และคณะ \cite{pascucci2021,pascucci2020preprint} นับเป็นกลุ่มแรกที่พัฒนาแอปพลิเคชันมือถือที่ใช้ AI สำหรับวิเคราะห์การดื้อยาปฏิชีวนะแบบอัตโนมัติ งานวิจัยดังกล่าวเป็นเครื่องพิสูจน์ว่า เพียงแค่กล้องสมาร์ทโฟนที่ทำงานร่วมกับ Deep Learning ก็สามารถประเมินค่า ZOI ได้อย่างแม่นยำในระดับที่ยอมรับได้ทางคลินิก ซึ่งมีประโยชน์อย่างยิ่งสำหรับพื้นที่ที่มีทรัพยากรจำกัด

\begin{figure}[H]
  \centering
  \safeimg[width=\textwidth]{figures/_review/related_work_existing_system.png}{Existing automated AST system}
  \caption{ตัวอย่างระบบวิเคราะห์ AST อัตโนมัติเชิงพาณิชย์ที่มีอยู่แล้ว แสดงผลการตรวจจับ
    ZOI พร้อมตารางผลการจำแนก S/I/R ที่ใช้เป็นเกณฑ์เปรียบเทียบ \cite{rajasekar2023}}
  \label{fig:related_existing}
\end{figure}
Rajasekar และคณะ \cite{rajasekar2022} ได้นำเสนอ Ab.ai ซึ่งเป็นเครื่องมือ AI เฉพาะทางที่ออกแบบมาเพื่อรายงาน Antibiograms แบบอัตโนมัติ โดยสามารถทำความแม่นยำได้สูงเมื่อทดสอบกับจาน Disk Diffusion ในสภาวะควบคุมของห้องปฏิบัติการ
Yu และคณะ \cite{yu2025} ได้เสนอแนวทางใหม่ในการระบุขอบเขตของ ZOI ด้วยวิธี Hue-contrast ที่ขับเคลื่อนด้วย AI ซึ่งถือเป็นทางเลือกหนึ่งนอกเหนือจากการใช้วิธีตีกรอบ Bounding Box แบบเดิม
Callebaut และคณะ \cite{callebaut2024} ได้ประเมินประสิทธิภาพของระบบ Radian\textregistered{} in-line carousel สำหรับการทดสอบ Disk Diffusion แบบอัตโนมัติ โดยใช้ผลการอ่านด้วยมือของผู้เชี่ยวชาญเป็นเกณฑ์อ้างอิง ซึ่งโครงงานนี้ได้นำเกณฑ์มาตรฐานดังกล่าวมาใช้เป็นเป้าหมายในการเปรียบเทียบเช่นกัน

Rajasekar และคณะ \cite{rajasekar2023} ยังได้ทำการทบทวนวรรณกรรมอย่างเป็นระบบ (Systematic Review) เกี่ยวกับวิธีการวิเคราะห์ภาพ AST ด้วย AI ที่เคยตีพิมพ์มาทั้งหมด และพบข้อสรุปที่น่าสนใจว่า โมเดล Deep Learning มีประสิทธิภาพเหนือกว่าเทคนิคประมวลผลภาพแบบคลาสสิก (เช่น Hough Transform, Active Contours และ Template Matching) อย่างสม่ำเสมอและมีนัยสำคัญ โดยเฉพาะเมื่อต้องเผชิญกับสถานการณ์จริงในคลินิกที่มีความซับซ้อน เช่น จานเพาะเชื้อที่มีแผ่นดิสก์ยาจำนวนมาก ZOI ที่ทับซ้อนกัน และลักษณะการเจริญเติบโตของแบคทีเรียที่ไม่สม่ำเสมอ การทบทวนนี้ได้ครอบคลุมตั้งแต่วิธีการในอดีตที่นำ CNN ไปประมวลผลเฉพาะส่วนของภาพแผ่นดิสก์ที่ถูกตัดมา (Cropped Image) ไปจนถึงระบบ Pipeline การตรวจจับแบบครบวงจร (End-to-end) ซึ่งมีลักษณะคล้ายคลึงกับสถาปัตยกรรมที่พัฒนาขึ้นในโครงงานนี้

หนึ่งในข้อค้นพบที่สำคัญที่สุดของ Rajasekar และคณะ คือ Deep Learning แสดงให้เห็นถึงประสิทธิภาพที่เพิ่มขึ้นอย่างก้าวกระโดดเมื่อต้องรับมือกับปัญหาโซนทับซ้อน (Overlapping Zones)---ซึ่งเป็นเงื่อนไขที่ทำลายสมมติฐานเรื่อง "ความเป็นวงกลมที่สมบูรณ์แบบและแยกขาดจากกัน" อันเป็นหัวใจหลักของวิธี Hough Transform---รวมถึงมีความทนทาน (Robustness) ต่อความแปรปรวนของสภาพแวดล้อมในการถ่ายภาพ เช่น แสงที่ไม่สม่ำเสมอ มุมกล้องที่บิดเบี้ยว และ Noise จากการบีบอัดไฟล์ภาพ \cite{rajasekar2023} ข้อมูลเหล่านี้เป็นเครื่องยืนยันความถูกต้องในการตัดสินใจออกแบบสถาปัตยกรรมสำหรับโครงงานนี้ นั่นคือ การเลือกใช้ Deep Learning เป็นแกนหลักในการตรวจจับวัตถุ และจัดวางให้วิธี Hough Transform เป็นเพียงเกณฑ์มาตรฐาน (Baseline) สำหรับการเปรียบเทียบ มากกว่าที่จะเป็นองค์ประกอบหลัก นอกจากนี้ งานวิจัยดังกล่าวยังชี้ให้เห็นว่า ข้อจำกัดสำคัญของระบบ Deep Learning ส่วนใหญ่มาจากชุดข้อมูล AST ในอดีตที่มีขนาดค่อนข้างเล็ก (มักจะมีไม่ถึง 500 ภาพ) ประเด็นนี้จึงกลายเป็นแรงผลักดันให้โครงงานนี้ทุ่มเทความสำคัญให้กับการขยายข้อมูล (Data Augmentation) และสร้างข้อมูลสังเคราะห์ (Synthetic Data) เพื่อขยายขนาดชุดข้อมูลสำหรับการฝึกสอนโมเดลให้มีมากกว่า 2,000 ภาพ

\subsection{มาตรฐานการประเมินสำหรับวิธี AST อัตโนมัติ (Evaluation Standards for Automated AST Methods)}

Humphries และคณะ \cite{humphries2018} ได้นำเสนอประเด็นสำคัญเชิงระเบียบวิธีเกี่ยวกับแนวทางที่เหมาะสมในการประเมินระบบการวัด AST อัตโนมัติ เพื่อเปรียบเทียบกับเกณฑ์มาตรฐานทองคำ (Gold Standard) ซึ่งก็คือการวัดด้วยการกะระยะสายตาหรือไม้บรรทัดโดยนักจุลชีววิทยาที่ผ่านการฝึกอบรม เอกสารแนวทางปฏิบัติที่เป็นเลิศ (Best Practice) ของพวกเขา ซึ่งจัดทำโดยคณะทำงานพัฒนาและกำหนดมาตรฐานวิธีปฏิบัติของ CLSI ได้ระบุเกณฑ์การประเมินสำคัญหลายประการที่โครงงานนี้นำมาประยุกต์ใช้ สิ่งที่สำคัญที่สุดสำหรับงานวิจัยนี้คือ พวกเขาได้กำหนดว่า ค่าความคลาดเคลื่อนสัมบูรณ์เฉลี่ย (MAE) ที่เกิดจากระบบอัตโนมัติเมื่อเทียบกับการวัดด้วยมือ จะต้องมีค่าน้อยกว่า 3\,mm จึงจะถือว่าอยู่ในระดับที่ ``ยอมรับได้ทางคลินิก'' (Clinically Acceptable)---กล่าวคือ หากระบบมีข้อผิดพลาดเฉลี่ยต่ำกว่าระดับนี้ โอกาสที่ระบบจะแปลผล S/I/R คลาดเคลื่อนไปจากความเป็นจริงสำหรับคู่ยาและเชื้อแบคทีเรียส่วนใหญ่ก็จะมีน้อยมากจนไม่มีนัยสำคัญ

Humphries และคณะ ยังเน้นย้ำถึงความจำเป็นที่จะต้องประเมินระบบอัตโนมัติบนชุดข้อมูลทดสอบที่มีความหลากหลาย ทั้งในแง่ของสายพันธุ์แบคทีเรีย ชนิดของกลุ่มยาปฏิชีวนะ และรูปแบบการจัดวางแผ่นดิสก์บนจานเพาะเชื้อ เนื่องจากข้อผิดพลาดในการวัดขนาดเพียงเล็กน้อยสำหรับคู่ยา-เชื้อคู่หนึ่ง อาจรุนแรงถึงขั้นเปลี่ยนผลการตัดสินใจทางคลินิกสำหรับคู่อื่นได้ \cite{humphries2018} ข้อเสนอแนะนี้ส่งผลโดยตรงต่อการออกแบบการประเมินผลของโครงงานนี้ ซึ่งจะรายงานค่า MAE ที่คำนวณจากชุดข้อมูลทดสอบจำนวน 235 ภาพ (ในเฟสการฝึกสอน) และเพิ่มเติมในชุดข้อมูลทดสอบอิสระ ที่ครอบคลุมความหลากหลายจากภาพจานเพาะเชื้อของจริง นอกจากนี้ การประเมินยังต้องแยกแยะระหว่างข้อผิดพลาดแบบสุ่ม (Random Errors) และความลำเอียงเชิงระบบ (Systematic Bias หรือทิศทางของข้อผิดพลาด) เนื่องจากแนวโน้มของระบบที่มักจะวัดขนาดโซนใหญ่เกินไป (Overestimation) หรือเล็กเกินไป (Underestimation) อย่างสม่ำเสมอ สามารถแก้ไขได้ด้วยการปรับตั้งค่าหลังการประมวลผล (Post-processing Calibration)---ซึ่งตรงกับวัตถุประสงค์ของขั้นตอนการปรับตั้งค่าชดเชย (Calibration Tuning) ที่ได้ดำเนินการไว้ในระบบนี้อย่างพอดี
